% GENERATED FILE. Edit canonical lessons, then run npm run generate:books.
% canonical-source-hash: fnv1a64:e1bb3352d2da620c
% canonical-lessons: LA-C172-sero, LA-C172-meto, LA-C172-venor, LA-C172-piscor, LA-C172-nego

\chapter{To sow, To reap, To hunt, To go fishing, To deny}
\label{ch:la-to-sow-to-reap-to-hunt-to-go-fishing-to-deny}
\hlchaptermodality{172}

\begin{chapteropening}
\textbf{By the end of this chapter:} \emph{I can say to sow, to reap, to hunt, to go fishing and to deny.}

Complete the last lesson of chapter 172: I can say to sow, to reap, to hunt, to go fishing and to deny.
\end{chapteropening}

\section[serō, serere]{serō, serere — to sow, to plant}

\label{lesson:LA-C172-sero}



\begin{quote}
Before the new one: say the Latin for to create, then the Latin for to believe.
\end{quote}

\subsection*{You'll want to know: serō, serere}
\textbf{serō, serere} — \textquotedblleft{}to sow, to plant\textquotedblright{}.

\begin{cousinweb}
\textbf{serō}, \textquotedblleft{}to sow\textquotedblright{}, is inherited from Proto-Italic \emph{*sizō}, \textquotedblleft{}I sow\textquotedblright{}, usually traced to Proto-Indo-European \emph{*si-sh1-}, \textquotedblleft{}to sow\textquotedblright{}, a doubled form of the root (EtymDB 2.1); \emph{sēmen}, \textquotedblleft{}seed\textquotedblright{}, has the same root. It is a different word from \emph{serō}, \textquotedblleft{}to join, to bind\textquotedblright{}.
\end{cousinweb}

\begin{grammarlens}[title={What you've built}]
One more word for this chapter.
\end{grammarlens}

\subsection*{Your turn}
\begin{itemize}
\raggedright
  \item \emph{Say it:} \emph{serō, serere}
  \item \emph{Say it:} \emph{serō, serere}, once more
  \item \emph{Say it:} say \emph{serō, serere}
  \item \emph{Recall:} say the Latin for to create, then the Latin for to believe, then say \emph{serō, serere} again
\end{itemize}

\begin{tcolorbox}[breakable,colback=teal!4,colframe=teal!35!black,title={Before you move on}]
What does serō, serere mean? (\textquotedblleft{}To sow, to plant\textquotedblright{}.) Say it once more.
\end{tcolorbox}
\section[metō, metere]{metō, metere — to reap, to harvest}

\label{lesson:LA-C172-meto}



\begin{quote}
Before the new one: say the Latin for to believe, then the Latin for to sow.
\end{quote}

\subsection*{You'll want to know: metō, metere}
\textbf{metō, metere} — \textquotedblleft{}to reap, to harvest\textquotedblright{}.

\begin{cousinweb}
\textbf{metō} is inherited, usually traced to Proto-Indo-European \emph{*h2meh1-}, \textquotedblleft{}to mow, to reap\textquotedblright{} (EtymDB 2.1); Greek \emph{amaō}, \textquotedblleft{}to reap\textquotedblright{}, and English \emph{mow} are cousins (Lewis \& Short).
\end{cousinweb}

\begin{grammarlens}[title={What you've built}]
One more word for this chapter.
\end{grammarlens}

\subsection*{Your turn}
\begin{itemize}
\raggedright
  \item \emph{Say it:} \emph{metō, metere}
  \item \emph{Say it:} \emph{metō, metere}, once more
  \item \emph{Say it:} say \emph{metō, metere}
  \item \emph{Recall:} say the Latin for to believe, then the Latin for to sow, then say \emph{metō, metere} again
\end{itemize}

\begin{tcolorbox}[breakable,colback=teal!4,colframe=teal!35!black,title={Before you move on}]
What does metō, metere mean? (\textquotedblleft{}To reap, to harvest\textquotedblright{}.) Say it once more.
\end{tcolorbox}
\section[vēnor, vēnārī]{vēnor, vēnārī — to hunt}

\label{lesson:LA-C172-venor}



\begin{quote}
Before the new one: say the Latin for to sow, then the Latin for to reap.
\end{quote}

\subsection*{You'll want to know: vēnor, vēnārī}
\textbf{vēnor, vēnārī} — \textquotedblleft{}to hunt\textquotedblright{}.

\begin{cousinweb}
\textbf{vēnor} has no sure origin; EtymDB 2.1 records proposed links with Proto-Indo-European \emph{*wenh1-}, \textquotedblleft{}to strive for, to desire\textquotedblright{}, and \emph{*weyh1-}, \textquotedblleft{}to pursue\textquotedblright{}. English \emph{venison}, \textquotedblleft{}hunted meat\textquotedblright{}, comes from it through French.
\end{cousinweb}

\begin{grammarlens}[title={What you've built}]
One more word for this chapter.
\end{grammarlens}

\subsection*{Your turn}
\begin{itemize}
\raggedright
  \item \emph{Say it:} \emph{vēnor, vēnārī}
  \item \emph{Say it:} \emph{vēnor, vēnārī}, once more
  \item \emph{Say it:} say \emph{vēnor, vēnārī}
  \item \emph{Recall:} say the Latin for to sow, then the Latin for to reap, then say \emph{vēnor, vēnārī} again
\end{itemize}

\begin{tcolorbox}[breakable,colback=teal!4,colframe=teal!35!black,title={Before you move on}]
What does vēnor, vēnārī mean? (\textquotedblleft{}To hunt\textquotedblright{}.) Say it once more.
\end{tcolorbox}
\section[piscor, piscārī]{piscor, piscārī — to go fishing}

\label{lesson:LA-C172-piscor}



\begin{quote}
Before the new one: say the Latin for to reap, then the Latin for to hunt.
\end{quote}

\subsection*{You'll want to know: piscor, piscārī}
\textbf{piscor, piscārī} — \textquotedblleft{}to go fishing\textquotedblright{}.

\begin{cousinweb}
\textbf{piscor} is the verb made from \emph{piscis}, \textquotedblleft{}a fish\textquotedblright{} (Lewis \& Short), which is usually traced to Proto-Indo-European \emph{*pisk-}, \textquotedblleft{}fish\textquotedblright{} (EtymDB 2.1).
\end{cousinweb}

\begin{grammarlens}[title={What you've built}]
One more word for this chapter.
\end{grammarlens}

\subsection*{Your turn}
\begin{itemize}
\raggedright
  \item \emph{Say it:} \emph{piscor, piscārī}
  \item \emph{Say it:} \emph{piscor, piscārī}, once more
  \item \emph{Say it:} say \emph{piscor, piscārī}
  \item \emph{Recall:} say the Latin for to reap, then the Latin for to hunt, then say \emph{piscor, piscārī} again
\end{itemize}

\begin{tcolorbox}[breakable,colback=teal!4,colframe=teal!35!black,title={Before you move on}]
What does piscor, piscārī mean? (\textquotedblleft{}To go fishing\textquotedblright{}.) Say it once more.
\end{tcolorbox}
\section[negō, negāre]{negō, negāre — to deny, to say no}

\label{lesson:LA-C172-nego}



\begin{quote}
Before the new one: say the Latin for to hunt, then the Latin for to go fishing.
\end{quote}

\subsection*{You'll want to know: negō, negāre}
\textbf{negō, negāre} — \textquotedblleft{}to deny, to say no\textquotedblright{}.

\begin{cousinweb}
\textbf{negō} is built on the negative \emph{ne}, \textquotedblleft{}not\textquotedblright{}, perhaps through \emph{nec}, \textquotedblleft{}and not\textquotedblright{}; how the verb was formed is uncertain. English \emph{negative} and \emph{deny} come from its family.
\end{cousinweb}

\begin{grammarlens}[title={What you've built}]
One more word for this chapter.
\end{grammarlens}

\subsection*{Your turn}
\begin{itemize}
\raggedright
  \item \emph{Say it:} \emph{negō, negāre}
  \item \emph{Say it:} \emph{negō, negāre}, once more
  \item \emph{Say it:} say \emph{negō, negāre}
  \item \emph{Recall:} say the Latin for to hunt, then the Latin for to go fishing, then say \emph{negō, negāre} again
\end{itemize}

\begin{tcolorbox}[breakable,colback=teal!4,colframe=teal!35!black,title={Before you move on}]
What does negō, negāre mean? (\textquotedblleft{}To deny, to say no\textquotedblright{}.) Say it once more.
\end{tcolorbox}
